\section{Scalar composition along a path and at its endpoint}\label{sec:temporal}
The preceding constructions encode one source instance in one scalar.
We now distinguish three ways of using independent instances: requesting
several path diagnostics, compiling their answers into one endpoint
coordinate, and merely revisiting a fixed hidden input.  Only the first two
supply a direct-sum conclusion.  Throughout this section the input is the
\emph{formulation}, with the full SQ interface of
Section~\ref{sec:realizations}; an exact oracle for a later iterate is not
supplied without charge.

\subsection{A direct sum for classical formulation queries}
Let $g$ be the promised two-function Forrelation predicate, with
$|\Phi|\leq\alpha=1/100$ in case zero and $\Phi\geq\beta=3/5$ in
case one.  Write $L_g=R_{1/3}(g)$ for its worst-case bounded-error
randomized query complexity.  The cyclic instances have
\[
 L_g=\Omega\!\left(\frac{q^{\ell/2}}{r\ell}\right),\qquad
 a_i^TM_i^{-1}e=h\Phi_i,\qquad
 \norm{M_i}=1,\quad\sigma_{\min}(M_i)=K^{-1},
\]
where $h=\Theta(\delta)$ and the one-block dimension is
$N_0=\Theta(Tq^\ell)$.  Here $T$ remains the clock length of
Section~\ref{sec:realizations}; the positive integer $B$ denotes the number
of independent blocks.

\begin{lemma}[Classical direct sum]\label{lem:trajectory-direct-sum}
Suppose every query to a formulation built from $B$ independent promised
inputs is simulated by $O(1)$ classical queries, each touching one named or
sampled input block.  All choices of sampled block have public
probabilities.  An output from which all $B$ promise bits can be recovered,
with joint success at least $2/3$, requires $\Omega(BL_g)$ formulation
queries.
\end{lemma}
\begin{proof}
Constant-error randomized query complexities are comparable by worst-case
amplification.  Minimax therefore gives a distribution $\mu$ for which
any randomized algorithm using fewer than $cL_g$ queries has average
success at most $9/16$.  Embed a target drawn from $\mu$ in a uniform
hidden block $I$, and fill the other blocks independently from $\mu$.
The simulator may sample and retain complete filler inputs for free: only
queries to the target count.  The assembled input has distribution $\mu^B$
and is independent of $I$.  Each of the $O(Q)$ simulated queries made by
a $Q$-query formulation algorithm thus touches the target with average
probability $1/B$.  This also holds for global SQ samples with nonuniform
public block weights.  The expected target cost is $O(Q/B)$.
Truncate after twelve times that expectation.  Markov's inequality loses
at most $1/12$ success, leaving average target success at least $7/12$.
Since $7/12>9/16$, the chosen distribution forces $Q/B=\Omega(L_g)$.
No amplification of a distributional average advantage is used.
\end{proof}

This elementary argument is a special case of randomized direct-sum
phenomena \citep{BlaisBrody2019}.  It does not apply to one coherent
query in superposition over all blocks.

\subsection{Successive scalar increments of a box central path}
Consider the single LP
\begin{equation}\label{eq:multiplex-lp}
 \max\ 2\sum_{i=1}^B\theta_i e^TM_ix_i,
 \qquad -\mathbf1\leq M_ix_i\leq\mathbf1,
 \qquad \theta_i=\Gamma^{-i}.
\end{equation}
A public three-sparse partial-sum chain defines the free coordinate
$w=\sum_i a_i^Tx_i$, as in \eqref{eq:accumulator}.  Define
\[
 \rho(z)=\frac{2z}{1+\sqrt{1+4z^2}},\qquad \rho(0)=0.
\]
Since $\rho/(1-\rho^2)=z$, the unique logarithmic central point is
$x_i(\eta)=\rho(\eta\theta_i)M_i^{-1}e$.  At
$\eta_j=\Gamma^{j+1/2}$, $0\leq j\leq B$, its $B$ increments, indexed
by $j=1,\ldots,B$, obey
\begin{equation}\label{eq:increment-kernel}
 D_j:=w(\eta_j)-w(\eta_{j-1})
   =h\sum_{i=1}^B d_{j-i}\Phi_i,
 \qquad d_m=\rho(\Gamma^{m+1/2})-\rho(\Gamma^{m-1/2}).
\end{equation}
The kernel is nonnegative and has infinite sum one.  The identities
$\rho(z)\leq z$ and
$1-\rho(z)=\rho(z)/(z(1+\rho(z)))\leq1/(2z)$ imply
\begin{equation}\label{eq:increment-leakage}
 d_0\geq1-\frac{3}{2\sqrt\Gamma},\qquad
 \left|D_j/h-d_0\Phi_j\right|\leq1-d_0=:\tau_\Gamma.
\end{equation}
These bounds follow by telescoping; restricting the sum to the actual
blocks only reduces its off-diagonal mass.

\begin{theorem}[Scalar-increment trajectory]\label{thm:lp-trajectory}
For $\Gamma=2^{20}$, returning all $B$ increments in
\eqref{eq:increment-kernel} to additive error $0.21h$, jointly with
probability $2/3$, requires $\Omega(BL_g)$ classical full-SQ queries.
The canonical coherent source interface permits $O(B)$ queries at this
accuracy, or $O(B\log(2B))$ queries at the stronger accuracy $h/100$.
The LP has dimension $\Theta(BN_0)$, row and column sparsity $O(q)$,
and selected multiplier range $\log(\eta_B/\eta_0)=B\log\Gamma$.
\end{theorem}
\begin{proof}
The low and high intervals for $D_j/h$ are separated by at least
$d_0(\beta-\alpha)-2\tau_\Gamma>0.586$.  Error $0.21$ therefore still
identifies the bit.  Each matrix query touches one block; the added
readout chain and objective coefficients are public.  Row and column
norms and sampling probabilities follow from public squared magnitudes.
The chain has $O(BT)$ rows and constant incidence, so it preserves the
claimed dimension and sparsity.  Lemma~\ref{lem:trajectory-direct-sum}
proves the lower bound.

For the numerical upper, estimate each $\Phi_j$ to error $0.005$, with
failure probability $1/(3B)$, and return $hd_0\widehat\Phi_j$.
Equation~\eqref{eq:increment-leakage} bounds the total error by
$0.005h+0.0015h<h/100$.  Constant-accuracy source estimation followed by
median amplification costs $O(\log(2B))$ queries per block.
For the looser accuracy, use the coherent Boolean recovery theorem of
\citet[Corollary~3]{BuhrmanEtAl2007}.  It recovers all $B$ bits jointly
using $O(B)$ controlled calls to their bounded-error constant-query
subroutines and their inverses.  Return zero for a low bit and
$0.8hd_0$ for a high bit.  Their errors are at most
$(0.01d_0+\tau_\Gamma)h<0.012h$ and
$(0.2d_0+\tau_\Gamma)h<0.202h$, respectively.
\end{proof}

The Boolean recovery theorem is essential for the $O(B)$ claim: individual
constant success does not imply joint success.  Nor does Boolean recovery
supply arbitrary real-valued estimates.  In particular, absolute checkpoint
values include sums of actual amplitudes.  Estimating each amplitude to
$O(1/B)$ and failure $O(1/B)$ gives an elementary
$O(B^2\log(2B))$ quantum upper for absolute checkpoint accuracy $O(h)$;
no matching linear upper for that stronger output is asserted.

The same lower bound applies to exactly feasible approximate centers with
barrier-subproblem gap $\tau_{\rm ctr}\leq c h^2/(BK)$, for a sufficiently small
constant $c$.  Indeed, the interval barrier Hessian is at least $2I$ in
$u_i=M_ix_i$ coordinates, so strong convexity gives
$\norm{\widehat u-u}\leq\sqrt{\tau_{\rm ctr}}$.  Since
$\norm{M_i^{-T}a_i}\leq K/R\leq\sqrt K$,
\[
 |\widehat w-w|\leq\sqrt{BK\tau_{\rm ctr}}.
\]
Differencing two such values preserves the decoding gap.  The accumulator
equality must hold exactly, or its residual must be included in this error
budget.  At the analytic center the reduced Hessian has condition $K^2$;
there is no uniform condition bound over the whole scaled trajectory.
The paired interval barrier has parameter one as established in
Section~\ref{sec:realizations}, giving parameter $BN_0$ for this product.

For the fixed-$k$ promise in Corollary~\ref{cor:fixed-k}, choose
$\sqrt\Gamma\geq12/(\beta_k-\alpha_k)$ and use error
$c_kh$, $c_k>0$ sufficiently small.  The same proof gives
\[
 \widetilde\Omega_k(BN_0^{1-1/k})
 =\widetilde\Omega_k(B^{1/k}\mathcal N^{1-1/k}),
 \qquad\mathcal N=\Theta(BN_0),
\]
and a numerical coherent upper $2^{O(k)}B\log(2B)$.  Even with the
positive high promise, its interval $[\beta_k,1]$ can be much wider than
the decoding gap; one Boolean representative need not meet that numerical
accuracy.

\subsection{Ordinary SOCP predictor decrements}
Let $g_i=M_i^{-T}a_i$, and use the public common norm
$G=\norm{g_i}$ from Lemma~\ref{lem:tilt-norm}.  Put
\[
 v_i=e+\frac{g_i}{2G},\quad C=\frac54,\quad
 \xi=\frac{h}{CG},\quad\norm{v_i}^2=C(1+\xi\Phi_i),
 \quad 0<\xi\leq\frac14.
\]
The last restriction is obtained by taking the public upper bound on
$\delta$ small enough.  On the product of $B$ unit balls consider
\begin{equation}\label{eq:multiplex-socp}
 \max\ 2\sum_i\theta_i\left(e^Tu_i+\frac{a_i^Tx_i}{2G}\right),
 \qquad u_i=M_ix_i,\quad\norm{u_i}\leq1,
 \quad\theta_i=\frac{\Gamma^{-i}}{\sqrt C}.
\end{equation}
The barrier is $-\sum_i\log(1-\norm{u_i}^2)$, with the conventional
Lorentz parameter $2B$. At its exact central point for $\eta_j=\Gamma^j$,
with $j=1,\ldots,B$,
propose the relative update $\eta_j\mapsto(1+\sigma)\eta_j$, where
$\sigma=\sigma_0/\sqrt B$ and $\sigma_0$ is a small constant.

\begin{theorem}[Predictor-decrement trajectory]\label{thm:socp-trajectory}
For $\Gamma\geq256$, the ordinary squared predictor decrements for
$j=1,\ldots,B$ satisfy
\begin{equation}\label{eq:decrement-trajectory}
 P_j:=\frac{\Lambda_j^2}{\sigma^2}
 =\sum_i f\!\left(\Gamma^{j-i}\sqrt{1+\xi\Phi_i}\right),
 \qquad f(z)=1-\frac1{\sqrt{1+4z^2}}.
\end{equation}
Returning all $B$ predictor decrements $\Lambda_j^2$ to error
$\sigma^2\xi/100$, with joint
success $2/3$, requires $\Omega(BL_g)$ classical full-SQ formulation
queries.  Their scale is
$\sigma^2\xi=\Theta(\delta/(B\sqrt K))$, and every predictor obeys
$\Lambda_j^2\leq\sigma_0^2$.
The structured coherent upper is $O(B\log(2B))$.  For
$\Gamma=2^{20}$, at the looser error $0.036\sigma^2\xi$, it is $O(B)$,
while the same classical lower bound holds.
\end{theorem}
\begin{proof}
Write $s_i=\norm{v_i}$.  The central point is
$u_i=\rho(\eta\theta_i s_i)v_i/s_i$.  At radius $r_i$, the radial
barrier-Hessian eigenvalue is $2(1+r_i^2)/(1-r_i^2)^2$.
The multiplier update leaves residual $2\sigma\eta\theta_i v_i$.
Its squared inverse-Hessian norm is
$2\sigma^2r_i^2/(1+r_i^2)=\sigma^2 f(\eta\theta_i s_i)$.
This proves \eqref{eq:decrement-trajectory} and the bound on its size.

Subtract the public baseline $P_j^0=\sum_i f(\Gamma^{j-i})$ and set
$\psi_z(x)=f(z\sqrt{1+\xi x})-f(z)$.  Direct differentiation gives
\begin{equation}\label{eq:decrement-sensitivity}
 \psi_z'(x)=\frac{2\xi z^2}{(1+4z^2(1+\xi x))^{3/2}},\qquad
 \frac{2\xi}{6^{3/2}}\leq\psi_1'(x)\leq\frac\xi4
 \quad (|x|\leq1).
\end{equation}
The other blocks contribute absolute total at most
\begin{align}
 T_\Gamma
 &\leq\xi\left(\frac{2}{3^{3/2}(\Gamma-1)}
                 +\frac{2}{\Gamma^2-1}\right)
 \leq\frac{0.4\xi}{\Gamma-1},\qquad\Gamma\geq256.
 \label{eq:socp-leakage}
\end{align}
For scales $z=\Gamma^m\geq\Gamma$, use
$\psi_z'\leq2\xi/(3^{3/2}z)$; for $z\leq\Gamma^{-1}$, use
$\psi_z'\leq2\xi z^2$, and sum both geometric series.
The low/high gap for $P_j-P_j^0$ is at least
\[
 \xi\left(\frac{2\beta}{6^{3/2}}-\frac\alpha4
                  -\frac{0.8}{\Gamma-1}\right)>0.075\xi.
\]
Lemma~\ref{lem:trajectory-direct-sum} proves the classical lower bound.
The two objective pieces in each block occupy disjoint $u$ and $x$
registers, with public squared norm
$4\theta_i^2(1+\norm{a_i}^2/(4G^2))$.  Thus objective SQ, as well as
equality-matrix SQ, has constant source-query simulation cost.

Estimating each active amplitude to error $0.02$ and failure $1/(3B)$,
then returning $P_j^0+\psi_1(\widehat\Phi_j)$, gives error at most
$0.005\xi+T_\Gamma<0.01\xi$.  For $\Gamma=2^{20}$, recover all
promise bits jointly and output the midpoint of
$P_j^0+\psi_1([-\alpha,\alpha])$ or
$P_j^0+\psi_1([\beta,1])$.  On the low interval,
$\psi_1'\leq2\xi/(5-0.01)^{3/2}<0.180\xi$, so the midpoint error is
less than $0.0018\xi$.  On the positive high interval,
$\psi_1'\leq2\xi/5^{3/2}$, so its midpoint error is at most
$0.4\xi/5^{3/2}<0.035778\xi$.
Adding \eqref{eq:socp-leakage} keeps both below $0.036\xi$.
The gap is still greater than $0.079\xi>2(0.036\xi)$.
\end{proof}

For fixed $k$, the same proof uses
$\Delta_k=2\beta_k/6^{3/2}-\alpha_k/4>0$, chooses
$\Gamma_k\geq\max\{256,1+1.6/\Delta_k\}$, and takes error
$c_k\sigma^2\xi$ with $c_k<\Delta_k/8$.
It gives the same $\widetilde\Omega_k(BN_0^{1-1/k})$ classical lower
and $2^{O(k)}B\log(2B)$ numerical coherent upper.
The formulation dimension is $\Theta(BN_0)$ and sparsity $O(q)$.
The reporting points span $(B-1)\log\Gamma$ in logarithmic multiplier.
Intermediate short steps may be inserted; the diagnostic at each reporting
point concerns the small update $\sigma$, not a jump of size $\Gamma$.
No full-trajectory Hessian condition bound is asserted.

\begin{corollary}[Cost of a dynamic block-norm interface]\label{cor:dynamic-norm}
For an interface constructed from the formulation, the total cost of setup
and answering all $B$ active cone-block norm requests
$\norm{u_j(\eta_j)}$ to additive error $\xi/50$, with joint success
at least $2/3$, is $\Omega(BL_g)$ classical queries.
For a precompiled interface whose subsequent requests each cost $O(1)$,
this bounds setup cost plus $O(B)$.
\end{corollary}
\begin{proof}
The norm equals $\rho(\sqrt{1+\xi\Phi_j})$.  In this argument
$\rho<2/3$ and
$\rho'(z)=(1-\rho^2)^2/(1+\rho^2)\geq25/117$.
Consequently its derivative in $\Phi_j$ is at least
$25\xi/(117\sqrt5)$.  The promise gap in the norm is greater than
$0.056\xi>2\xi/50$, so Lemma~\ref{lem:trajectory-direct-sum} applies.
\end{proof}

Free exact iterate SQ with conditional block norms would expose this
information.  Theorem~\ref{thm:socp-trajectory} therefore cannot be
reinterpreted as a lower bound with that stronger input already supplied.
The corollary charges the construction of precisely that interface.  A
normalized quantum state with its norm withheld is a different output.

\subsection{Compiling promised overlaps into one bounded optimizer coordinate}
A sum of promised amplitudes need not determine a sum of promise bits:
intervals corresponding to different bit patterns can overlap.  The
following constant-size LP removes that obstacle.

\begin{lemma}[LP threshold compiler]\label{lem:threshold-lp}
Suppose $\varphi\in[-1,1]$ satisfies $|\varphi|\leq\alpha$ or
$\varphi\geq\beta$, with $0\leq\alpha<\beta\leq1$.
Put $c=(\alpha+\beta)/2$, $\Delta=(\beta-\alpha)/2$,
$J=1/\Delta$, $M_0=2J$, and $C_0=4J$.
The LP
\begin{equation}\label{eq:threshold-lp}
 \max\ C_0z-M_0r+t,\qquad
 0\leq z\leq1,\quad0\leq r\leq2,\quad
 0\leq t\leq z,\quad t\leq Jr,\quad r\geq z(\varphi-c)
\end{equation}
has a unique optimizer with $(z,t)=(1,g)$, where $g$ is the promise bit.
Every feasible point has objective gap at least $|t-g|$.
\end{lemma}
\begin{proof}
For a low input the last constraint is redundant.  For fixed $z$, the
best $t$ is $\min\{z,Jr\}$ and the penalty coefficient $M_0=2J$
uniquely sets $r=t=0$.  Then $z=1$ is unique.  Moreover the gap is
$C_0(1-z)+M_0r-t\geq t$.
For a high input set $d=\varphi-c\geq1/J$.  Feasibility forces
$r\geq dz$, so the unique best choices are $r=dz,t=z$.
The remaining coefficient $C_0+1-M_0d$ is positive since $d<1$.
Thus $z=t=1$ uniquely.  For every feasible point,
$f\leq(C_0-M_0d)z+t$ and $C_0-M_0d>0$, giving
$f^*-f\geq1-t$.
\end{proof}

For each cyclic block, introduce $M_iy_i=z_ie$ and a public sparse
readout chain defining $q_i=a_i^Ty_i/h$.  These linear equalities enforce
$q_i=z_i\Phi_i$.  Replace the last constraint of
\eqref{eq:threshold-lp} by $r_i\geq q_i-cz_i$.
There is no product of two optimization variables: $\Phi_i$ is implicit
input data, and its product with $z_i$ is implemented by linear equations.
Give block $i$ the objective weight $\theta_i>0$.
Set $p_1=t_1$ and, for $2\leq i\leq B$, impose
\begin{equation}\label{eq:xor-facets}
 \begin{split}
 p_i&\leq p_{i-1}+t_i,\qquad p_i\geq p_{i-1}-t_i,\\
 p_i&\geq t_i-p_{i-1},\qquad p_i\leq2-p_{i-1}-t_i.
 \end{split}
\end{equation}
These are the convex hull of Boolean XOR triples.  Their projection onto
$(p_{i-1},t_i)$ is the whole unit square, and Boolean inputs uniquely
force $p_i=p_{i-1}\mathbin\oplus t_i$.

\begin{theorem}[A bounded endpoint scalar]\label{thm:threshold-xor}
The compiled LP has a unique optimizer and
$P^*:=p_B^*=g_1\mathbin\oplus\cdots\mathbin\oplus g_B$.
Its dimension is $\Theta(BN_0)$, row and column sparsity $O(q)$,
and its displayed logarithmic barrier has certified parameter $12B-4$.
Estimating $P^*$ to additive error less than $1/3$ requires
$\Omega(BL_g)$ classical full-SQ queries and has coherent source query
complexity $\Theta(B)$ for two-function Forrelation.  For fixed $k$,
the classical lower is $\widetilde\Omega_k(BN_0^{1-1/k})$ and the
coherent complexity is $\Theta_k(B)$.
Every exactly feasible point with objective gap $\tau$ satisfies
\begin{equation}\label{eq:xor-gap}
 |p_B-P^*|\leq\frac{\tau}{\min_i\theta_i}.
\end{equation}
\end{theorem}
\begin{proof}
The XOR constraints leave the projection onto all $t_i$ unchanged.
Lemma~\ref{lem:threshold-lp} fixes each block's optimizer uniquely;
then the XOR variables, inverse histories and accumulators are also unique.
There are eight scalar inequalities per threshold block and four per later
XOR gate.  Strict relative feasibility follows by taking $z_i=3/4$,
$r_i=3/2$, $t_i=p_i=1/2$, with the equality variables determined
accordingly.  The resulting scalar product log barrier has parameter at
most the number $12B-4$ of its logarithms, independent of $N_0$.

For $B=1$, the randomized lower is the one-block lower bound immediately.
For $B\geq2$, it follows from the strong XOR theorem of
\citet[Theorem~1.1]{BrodyEtAl2023}:
$\overline R_\epsilon(\mathsf{XOR}_B\circ g^B)
=\Theta(B\overline R_{\epsilon/B}(g))$ for constant outer error,
where the bar denotes worst-input expected query cost.
Expected constant-error and worst-case constant-error costs are comparable:
truncate after a constant multiple of the worst-input expectation and
amplify the resulting worst-input success.  Hence the theorem implies
$\Omega(BR_{1/3}(g))$ also for worst-case queries.
For the coherent upper, jointly recover all promise bits using
\citet[Corollary~3]{BuhrmanEtAl2007} and take their XOR.  This costs
$O(B)$ or $2^{O(k)}B$ source queries.  For the lower, restrict every
block to two fixed inputs having opposite promise bits.  One source query
on this restriction is simulated by $O(1)$ queries to the selecting outer
bit.  Thus an algorithm computes parity of $B$ independent bits, which
needs $\Omega(B)$ quantum queries \citep{BealsEtAl2001}.

Finally, Lemma~\ref{lem:threshold-lp} gives
$\tau\geq\sum_i\theta_i|t_i-g_i|$.  The XOR facets imply
$|p_i-p_{i-1}|\leq t_i$ and
$|p_i-(1-p_{i-1})|\leq1-t_i$.
Induction, using the appropriate inequality for $g_i=0$ or $1$, yields
$|p_i-(g_1\mathbin\oplus\cdots\mathbin\oplus g_i)|
\leq\sum_{j\leq i}|t_j-g_j|$, proving \eqref{eq:xor-gap}.
\end{proof}

All new coefficients and supports are public.  An old row or column gains
only constantly many public entries, so its new SQ law is a public mixture
of the old law and those entries.  Readout chains keep incidence constant.
Thus full formulation SQ has constant source-query overhead.  Conversely,
hidden signs remain at known matrix entries, so the source upper has
matched formulation access.  The factor $1/h$ is public; no condition-number
bound for the entire compiled equality system is claimed.

With $\theta_i=1$, constant objective gap $\tau<1/3$ suffices for the
same coordinate lower bound.  With $\theta_i=\Gamma^{-i}$, the sufficient
gap is $\tau<\Gamma^{-B}/3$.  Before adding XOR facets the separated
threshold blocks have central problems at multipliers $\eta\theta_i$;
a public summation readout $\sum_i t_i$ does not affect this property.
The XOR facets couple the barrier, so this rescaling identity is not
asserted for the bounded-output formulation.  Endpoint hardness does not
establish an unavoidable charge at each of $B$ path scales.

The compiler tolerates finite public readout precision.  If
$b_i=a_i/h$ is replaced by $\widetilde b_i$ with
$\norm{\widetilde b_i-b_i}\leq\Delta/(10R)$, then
$\norm{y_i}\leq R$ makes the readout error at most $\Delta/10$.
Use the narrower effective half-gap $9\Delta/10$ when choosing $J,M_0,C_0$.
This preserves exact optimizer bits: although a rounded amplitude can
slightly exceed one, its high-branch value $d=\widetilde\Phi_i-c$ is
still below one because $c\geq\Delta$, so the proof of
Lemma~\ref{lem:threshold-lp} applies with the narrower gap.
Keeping the old threshold constants
need not preserve saturation at the high boundary.  This statement rounds
only public readouts; rounding $M_i$ requires a separate resolvent budget.

\subsection{A parity endpoint and the meaning of temporal hardness}
A simpler sparse family illustrates the distinction without an inner
quantum separation.  Let $\sigma_{i,k}\in\{-1,1\}$,
$1\leq i\leq B$, $1\leq k\leq L$, and
$h_i=\prod_k\sigma_{i,k}$.  Introduce boxed activations
$-1\leq a_i\leq1$ and chains
\[
 x_{i,0}=a_i,\qquad x_{i,k}=\sigma_{i,k}x_{i,k-1}.
\]
A public partial-sum chain defines $W=\sum_i x_{i,L}$.
Order all $BL$ signs as $\tau_1,\ldots,\tau_{BL}$, and add the carrier
$r_0=a_B$, $r_m=\tau_m r_{m-1}$, $w=r_{BL}$.
Maximize $2\sum_i\Gamma^{-i}a_i$, with $\Gamma=256$.
All equalities have at most three nonzeros per row and constant column
incidence.  There are $2B$ inequalities and $\Theta(BL)$ variables.
The optimizer is unique, and direct elimination gives
\begin{equation}\label{eq:parity-path}
 W(\eta)=\sum_i h_i\rho(\eta\Gamma^{-i}),\qquad
 w(\eta)=H\rho(\eta\Gamma^{-B}),\qquad
 w^*=H:=\prod_i h_i.
\end{equation}
At $\eta_j=\Gamma^{j+1/2}$, the increment kernel
\eqref{eq:increment-kernel} gives
$h_j(W(\eta_j)-W(\eta_{j-1}))\geq2d_0-1\geq13/16$.
Thus the path has $B$ separated constant-margin transitions.
At its final checkpoint,
$|w(\eta_B)-H|\leq1/32$, so additive error $1/4$ in this one bounded
coordinate computes parity of all $BL$ signs.  Both its quantum and
classical query complexities are $\Theta(BL)$; reading all signs proves
the upper bound and the parity lower proves the lower.
All nonzero magnitudes, norms and sampling laws are public, and each
queried matrix entry uses at most one source sign.  Duplicate occurrences
refer to the same fixed input bit.  The barrier parameter $2B$ and public
$O(B)$-bit objective scales do not alter this query accounting.

Every feasible point of objective gap $\tau$ has
$|w-H|=1-a_B\leq\tau/(2\Gamma^{-B})$.  Equal objective weights remove
this exponentially small gap requirement but also remove the scale
separation.  At $\eta_{\mathrm f}=16B\Gamma^B$,
$|W(\eta_{\mathrm f})-\sum_i h_i|\leq1/32$; additive error $1/4$ in
this unnormalized sum also computes parity by rounding.  Normalizing by
$B$ changes the required accuracy to $O(1/B)$.
More generally, if a block bit has an exact coherent evaluator of cost
$q_0$, amplitude estimation gives its mean to error $\epsilon$ using
$O(q_0\min\{B,1/\epsilon\})$ queries
\citep{BrassardEtAl2002}.  A mean at constant accuracy therefore need not
retain a factor $B$, whereas the bounded parity carrier does.

The query lower bounds in this section concern specified final outputs or
joint transcripts.  In a static-oracle model an algorithm required only to
return a final relation may compute it without generating an intermediate
iterate.  A geometric transition is not a new oracle input.  A genuine
chronological charge needs an online-output, fresh-input, or memory
restriction, or a composition theorem for the required final function.
Our parity and threshold constructions use the last option; they do not
prove per-iteration lower bounds.

The Forrelation, direct-sum, robust recovery and XOR theorems used here are
prior work.  Optimization scalar encodings of Boolean composition also
appear in \citet{ApersGribling2026}.  Our contribution in this section is
the explicit sparse realization and its quantitative output contracts:
the increment leakage kernel, the ordinary predictor-decrement sensitivity
kernel, and the threshold compiler with objective-gap stability.  To the
best of our knowledge, these particular access-preserving central-path
realizations have not appeared in the cited literature; no priority claim
is made for their individual query-complexity ingredients.
